% ======================================================================
% Ontology of Continua -- Core 1.3.3
% Cross-K module: crossk_k2_k3.tex
% Cross-level relations between K2 and K3
% ======================================================================

\subsubsection{Межуровневая структура K2–K3}
\label{sec:crossk-k2-k3}

Переход $K_2 \rightarrow K_3$ описывает возникновение химического
континуума из физического континуума, уже обладающего геометрией, полями и
динамическими потоками. Если $K_2$ содержит физические степени свободы
(пространственные оси, энергетические потенциалы, градиенты полей, структуру
перколяции), то $K_3$ вводит:
\begin{itemize}
  \item пространства химического состава,
  \item энергии связи и молекулярные потенциалы,
  \item пороговые условия реакций и стехиометрические ограничения,
  \item устойчивые кластеры и сети связей,
  \item новые классы потоков (реакционные, диффузионно-реакционные,
        перенос заряда),
  \item расширенную область $\Omega(K_3)$, структурированную химической
        связностью.
\end{itemize}

Переход K₂→K₃ — это первый существенный скачок структурной сложности:
физические поля становятся основой химической организации.

% ----------------------------------------------------------------------
\subsubsection{1. Задействованные уровни и их роли}

\paragraph{K2 (физический континуум).}
$K_2$ предоставляет:
\begin{itemize}
  \item многомерные оси (геометрические координаты),
  \item энергетические потенциалы $P_2$,
  \item потоки $J_2$ (транспорт, диффузия, распространение поля),
  \item структуру перколяции связности,
  \item причинный порядок событий.
\end{itemize}

\paragraph{K3 (химический континуум).}
$K_3$ строится на основе $K_2$ и вводит:
\begin{itemize}
  \item новые оси, описывающие химические виды, состав и заряд,
  \item потенциалы, связанные с энергией связи и ландшафтом реакций,
  \item пороги $\Theta_3$ для реализуемости реакций и молекулярной
        устойчивости,
  \item потоки $J_3$, включающие реакции, диссоциацию и перенос заряда,
  \item устойчивые кластеры и реакционные циклы, задающие химические
        маршруты.
\end{itemize}

$K_3$ — первый континуум, поддерживающий шаблонизацию, катализ и
информационно-подобную персистентность через молекулярные состояния.

% ----------------------------------------------------------------------
\subsubsection{2. Совместные и наследуемые оси и пороги}

\paragraph{Оси.}
Пространственные оси $K_2$ полностью наследуются:
\[
A_{\mathrm{geom}}(K_2) \subset A(K_3).
\]
$K_3$ добавляет химические оси, такие как:
\begin{itemize}
  \item состав/стехиометрия,
  \item зарядовое состояние,
  \item конфигурация связей,
  \item координаты реакции.
\end{itemize}

\paragraph{Пороги.}
Химические пороги $\Theta_3$ расширяют физические пороги $\Theta_2$.
Ключевые унаследованные ограничения:
\begin{itemize}
  \item энергетические пороги: недостаточный $P_2$ не допускает образования
        связей;
  \item пороги перколяции: связность областей взаимодействия должна
        превысить критическое значение $p_c$;
  \item пороги градиента: требуется достаточный поток или частота столкновений
        для активации реакции.
\end{itemize}

Если физический порог нарушен, химическая структура не может возникнуть:
\[
\Theta_2^{\mathrm{death}} \Rightarrow \Omega(K_3)=\varnothing.
\]

% ----------------------------------------------------------------------
\subsubsection{3. Межуровневые потоки, циклы и напряжения}

\paragraph{Потоки.}
Потоки в $K_3$ построены на $J_2$, но теперь включают:
\begin{itemize}
  \item реакционные потоки по химическим координатам,
  \item диффузионно-реакционное сопряжение,
  \item перенос заряда и красные (redox) потоки,
  \item потоки между потенциальными ямами в ландшафте реакций.
\end{itemize}

Формально:
\[
J_3 = (J_2,\ J_{\mathrm{react}},\ J_{\mathrm{charge}},\ J_{\mathrm{bond}}).
\]

\paragraph{Циклы.}
Химические циклы имеют более высокую размерность, чем физические контуры.
Они представляют:
\[
C_3 : \text{реактивная последовательность } A \rightarrow B \rightarrow C \rightarrow A.
\]
Химические циклы требуют транспорта $K_2$ и реакционных траекторий $K_3$.

Существование устойчивого цикла в $K_3$ определяется:
\[
\oint_{\gamma} J_3 \cdot dA > 0
\quad\text{and}\quad
\text{all intermediate states satisfy }\Theta_3.
\]

\paragraph{Напряжение.}
Межуровневое напряжение $T_{2,3}$ измеряет совместимость между физическими
полями и химическими ограничениями:
\[
T_{2,3} = f(P_3 - P_2,\ \Theta_3 - \Theta_2,\ A_{\mathrm{chem}}).
\]
Если $T_{2,3}$ превышает порог, химический континуум схлопывается обратно в
чисто физический режим (без устойчивых молекул или реакций).

% ----------------------------------------------------------------------
\subsubsection{4. Условия рождения и исчезновения между уровнями}

\paragraph{Рождение \texorpdfstring{$K_3$}{K_3}.}
Химическая структура возникает, когда:
\[
T_2 > \Theta_{2,\mathrm{dim}}
\quad \text{и} \quad
A_{\mathrm{chem}} \in M_3 \setminus A(K_2).
\]
Это соответствует появлению новых степеней свободы, связанных с
энергиями связи и составом.

Пространство состояний расширяется:
\[
\Omega(K_3) = \Omega(K_2) \cup \Delta\Omega_{\mathrm{chem}}.
\]

\paragraph{Распространение гибели.}
Если $\Theta_2$ не выполняется, химическая организация невозможна.
Если $\Theta_3$ не выполняется (например, из-за недостаточной энергии связи,
несовместимых полей или чрезмерного теплового шума), разрушается только $K_3$,
в то время как $K_2$ остаётся не затронутым.

% ----------------------------------------------------------------------
\subsubsection{5. Концептуальные примеры}

\paragraph{Перколяция, обеспечивающая химическую связность.}
Лишь при перколяции областей физического взаимодействия ($p>p_c$) могут
формироваться химические кластеры; иначе молекулы не стабилизируются.

\paragraph{Ландшафты реакций.}
Переход вводит новые потенциалы:
\[
P_3 = P_2 + E_{\mathrm{bond}} + E_{\mathrm{charge}},
\]
которые создают потенциальные ямы и барьеры активации, задающие траектории реакций.

\paragraph{Операторная картина (создание/аннигиляция кластеров).}
Формирование или диссоциация минимального кластера $q_{\mathrm{min}}$ можно
выразить через операторы:
\[
a^{\dagger} |k\rangle = \sqrt{f(k)+1}\,|k+k_{\mathrm{unit}}\rangle,
\qquad
a |k\rangle = \sqrt{f(k)}\,|k-k_{\mathrm{unit}}\rangle,
\]
отражая квантизацию изменения связности внутри $\Omega(K_3)$.

Эти примеры выражают суть перехода K2→K3 в ограниченной форме:
устойчивые кластеры, энергии связи и реакционные маршруты возникают из
физических полей и геометрии.
